\section{交错多重线性形式的一个推广}

设 \(\mathrm{V}\) 是体 \(\mathrm{D}\) 上的一个右向量空间，维数有限，为 \(n \geqslant 0\)。用 \(\mathrm{B(V)}\) 表示 \(\mathrm{V}\) 的基所成的集合，用 \(\Omega(\mathrm{V})\) 表示满足下列两个条件的映射 \(\omega\)（从 \(\mathrm{B(V)}\) 到 \(\mathrm{D}_{\mathrm{ab}}^{*}\)）所成的集合：

\begin{enumerate}
\def\labelenumi{\alph{enumi})}
\tightlist
\item
  若 \(\lambda_1, \ldots, \lambda_n\) 是 \(D^*\) 中的元素，则我们有
\end{enumerate}

\[
\omega (v _ {1} \lambda_ {1}, \ldots , v _ {n} \lambda_ {n}) = \pi (\lambda_ {1} \dots \lambda_ {n}) \omega (v _ {1}, \ldots , v _ {n})\tag{1}
\]

对 V 的每个基 \((v_{1},\ldots ,v_{n})\) 成立。

\begin{enumerate}
\def\labelenumi{\alph{enumi})}
\setcounter{enumi}{1}
\tightlist
\item
  若 \(i\) 与 \(j\) 是区间 \([1, n]\) 中两个不同的整数，则我们有
\end{enumerate}

\[
\omega (v _ {1}, \ldots , v _ {i - 1}, v _ {i} + v _ {j}, v _ {i + 1}, \ldots , v _ {n}) = \omega (v _ {1}, \ldots , v _ {n})\tag{2}
\]

对 V 的每个基 \((v_{1},\ldots ,v_{n})\) 成立。

设 \(\omega\) 是 \(\Omega(\mathrm{V})\) 的一个元素。设 \((v_1, \ldots, v_n)\) 是 \(\mathrm{V}\) 的一个基，\(i\) 是区间 \([1, n-1]\) 中的一个整数；由性质 b)，我们有

\[
\begin{array}{r l} & {\omega (v _ {1}, \ldots , v _ {i}, v _ {i + 1}, \ldots , v _ {n}) = \omega (v _ {1}, \ldots , v _ {i} + v _ {i + 1}, v _ {i + 1}, \ldots , v _ {n})} \\ & {\qquad = \omega (v _ {1}, \ldots , v _ {i} + v _ {i + 1}, v _ {i + 1} - (v _ {i} + v _ {i + 1}), \ldots , v _ {n})} \\ & {\qquad = \omega (v _ {1}, \ldots , v _ {i} + v _ {i + 1}, - v _ {i}, \ldots , v _ {n})} \\ & {\qquad = \omega (v _ {1}, \ldots , (v _ {i} + v _ {i + 1}) - v _ {i}, - v _ {i}, \ldots , v _ {n}),} \end{array}
\]

从而

\[
\omega (v _ {1}, \dots , v _ {i}, v _ {i + 1}, \dots , v _ {n}) = \pi (- 1) \omega (v _ {1}, \dots , v _ {i + 1}, v _ {i}, \dots , v _ {n}).\tag{3}
\]

由于对称群 \(\mathfrak{S}_n\) 由区间 \([1,n]\) 中两个相邻元素的对换生成（I, §5, No.~7, 第 63 页，命题 9），公式 (3) 可推广为

\[
\omega (v _ {\sigma (1)}, \ldots , v _ {\sigma (n)}) = \pi (\varepsilon (\sigma)) \omega (v _ {1}, \ldots , v _ {n}),\tag{4}
\]

其中 \(\sigma\) 属于 \(\mathfrak{S}_n\)，而 \(\varepsilon (\sigma)\) 是它的符号差（I, §5, No.~7, 第 64 页）。

公式 (2) 可如下推广。首先，对每个 \(\lambda\)（在 \(\mathrm{D}^*\) 中），我们有

\[
\begin{array}{r l} & {\omega (v _ {1}, \ldots , v _ {n}) = \pi (\lambda) \omega (v _ {1}, \ldots , v _ {i} \lambda^ {- 1}, \ldots , v _ {n})} \\ & {\qquad = \pi (\lambda) \omega (v _ {1}, \ldots , v _ {i} \lambda^ {- 1} + v _ {j}, \ldots , v _ {n})} \\ & {\qquad = \omega (v _ {1}, \ldots , (v _ {i} \lambda^ {- 1} + v _ {j}) \lambda , \ldots , v _ {n}),} \end{array}
\]

也就是说，

\[
\omega (v _ {1}, \dots , v _ {n}) = \omega (v _ {1}, \dots , v _ {i} + v _ {j} \lambda , \dots , v _ {n}).\tag{5}
\]

由一个直接的归纳法，我们推出

\[
\omega (v _ {1}, \ldots , v _ {n}) = \omega (v _ {1}, \ldots , v _ {i} + w, \ldots , v _ {n})\tag{6}
\]

对每个线性组合 \(w = \sum_{j \neq i} v_j \lambda_j\)（诸向量 \(v_j\)，其中 j 在区间 \([1, n]\) 中且异于 i）成立。

